\documentclass[10pt,a4paper]{amsart}
\usepackage[margin=19mm]{geometry}
\usepackage{amsmath,amssymb,mathtools}
\usepackage[scr=rsfs,cal=euler]{mathalfa}
\usepackage{xfrac}
\usepackage[unicode,hidelinks,bookmarks=false]{hyperref}
\newcommand{\ie}{i.e.\ }
\newcommand{\cf}{cf.\ }
\newcommand{\resp}{respectively\ }
\newcommand{\Subsection}{\subsection}
\providecommand{\nobreakdash}{\nobreak\hbox{-}}
\setlength{\emergencystretch}{2em}
\multlinegap0pt
\usepackage{mathtools}
\usepackage{amssymb}
\usepackage{bm}
\usepackage[shortlabels]{enumitem}
\newtheorem{lem}{Lemma}
\newcommand{\GF}{\textup{GF}}
\newcommand{\bpd}{\overline{\partial}}
\newcommand{\ceq}{\coloneqq} %using mathtool package
\newcommand{\clE}{\mathcal{E}}
\newcommand{\ep}{\epsilon}
\newcommand{\lto}{\longrightarrow}
\title{BV construction of SUSY vertex algebras from SUSY factorization algebras}
\author{Shintarou Yanagida}
\date{Source: arXiv:2606.05706v1 (CC BY 4.0)}
\begin{document}
\maketitle

\noindent\emph{Source and licence.} BV construction of SUSY vertex algebras from SUSY factorization algebras. Version arXiv:2606.05706v1 (CC BY 4.0), distributed by arXiv under the Creative Commons Attribution 4.0 International licence, http://creativecommons.org/licenses/by/4.0/, as recorded in the arXiv licence metadata for this exact version and shown on the abstract page https://arxiv.org/abs/2606.05706v1. Full licence notice: \texttt{licenses/arxiv-2606.05706-CC-BY-4.0.txt}.\par\smallskip

\noindent\textit{Editorial scope.} Counted unit: the article's lem eq:2:hot (source0.tex lines 1802--1811). The article's complete original proof of it is delivered with it (source0.tex lines 1813--1848, one \texttt{proof} environment of 36 lines). No mathematics is written, repaired or re-derived: the statement and the proof are the article's own lines, in the article's own notation, and no retained line is substituted or re-flowed.  Retained uncounted as the source-local context the proof needs: lines 1798--1800.  Nothing was omitted from any retained span, so the package carries no omission record.  No retained line contains a citation, so the wrapper carries no bibliography and no bibliography entry.  No retained line contains a figure environment, an image-inclusion command or drawing code, so no image is delivered and the package declares no asset dependency.  Editorial formatting only, in addition: an AMS \texttt{amsart} preamble that restates the article's own declarations and macros used by the retained lines, namely the environments \texttt{lem} on separate counters, together with the standard packages those lines need.  The retained environments print in this document as Lemma 1 in order of appearance, whereas the article prints the same items with the numbering of its own full document.  The complete native source of the article as distributed in the arXiv e-print for this exact version is bundled unchanged as \texttt{sources/202135/source0.tex}.
\begin{align}\label{eq:2:PeL}
 P_{\ep<L} \ceq \int_\ep^L Q^{\GF} e^{-tH}\,dt \quad (0<\ep<L<\infty).
\end{align}

\begin{lem} %\label[lem]{lem:2:propagator}
For each $0<\ep<L<\infty$, the operator $P_{\ep<L}$ is 
a well-defined smoothing operator of cohomological degree $-1$. 
Moreover, it satisfies the homotopy identity
\begin{align}\label{eq:2:hot}
 [Q,P_{\ep<L}] = e^{-\ep H}-e^{-LH}.
\end{align}
In particular, $P_{\ep<L}$ is a heat-kernel regularized homotopy inverse 
for the differential $Q=\bpd$.
\end{lem}

\begin{proof}
This lemma is more or less a standard one, 
but we give a proof for completeness.
By construction, both $Q=\bpd$ and $Q^{\GF}=\bpd^{\,*}$ 
act only on the Dolbeault part of the space of fields, and do not affect
either the odd variable $\zeta$ or the coefficients in $V,V^\vee$. 
Hence
\[
 H = [Q,Q^{\GF}] = \bpd\,\bpd^{\,*} + \bpd^{\,*}\bpd
\]
is the Dolbeault Laplacian acting coefficient-wise on $\clE(U)$.

For each $t>0$, the heat operator $e^{-tH}$ is smoothing, 
and therefore $Q^{\GF}e^{-tH}$ is again smoothing. 
Hence the integral
$P_{\ep<L} = \int_\ep^L Q^{\GF}e^{-tH}\,dt$ in \eqref{eq:2:PeL} 
converges in the topology of smoothing operators, and defines a continuous map
\[
 P_{\ep<L}\colon \clE_c(U) \lto \clE(U)[-1].
\]
Its cohomological degree is $-1$, 
since $Q^{\GF}$ has degree $-1$ and $e^{-tH}$ has degree $0$.

It remains to prove \eqref{eq:2:hot}. 
Since $Q$ commutes with $H=[Q,Q^{\GF}]$, 
it also commutes with the heat operator $e^{-tH}$. 
Therefore
$[Q,Q^{\GF}e^{-tH}] = [Q,Q^{\GF}]\,e^{-tH} = He^{-tH}$.
Integrating it from $\ep$ to $L$, we obtain
\[
 [Q,P_{\ep<L}] = \int_\ep^L He^{-tH}\,dt =
 -\int_\ep^L \frac{d}{dt}e^{-tH}\,dt =
 e^{-\ep H}-e^{-LH},
\]
where we used $He^{-tH}=-\frac{d}{dt}e^{-tH}$ in the second equality.
\end{proof}

\end{document}
